\documentclass[14pt,a4paper]{extarticle}
\usepackage[T2A]{fontenc}
\usepackage[utf8]{inputenc}
\usepackage[russian]{babel}

% Стандартные академические поля ГОСТ (левое 30мм для подшивки, правое 10мм)
\usepackage[left=30mm,right=10mm,top=20mm,bottom=20mm]{geometry}
\usepackage{setspace}
\usepackage{indentfirst}
\usepackage{titlesec}

% Настройка полуторного интервала, как в реферате
\onehalfspacing

% Настройка красивых и компактных заголовков разделов резюме
\titleformat{\section}{\normalfont\large\bfseries}{\thesection.}{0.5em}{}
\titlespacing{\section}{\parindent}{1ex plus 0.5ex minus 0.2ex}{1ex plus 0.2ex}

\begin{document}

% В верхнем правом углу — ФИО студента и группа
\begin{flushright}
    \begin{spacing}{1.0}
        {\large Студент: Кузнецов М.А.}\\
        {\large Группа: ИУ7-42Б}
    \end{spacing}
\end{flushright}

\vspace{0.5cm}

% По центру — название реферата
\begin{center}
    \begin{spacing}{1.2}
        {\bf\Large РЕЗЮМЕ РЕФЕРАТА} \\
        \vspace{0.2cm}
        {\bf\large на тему: <<Истоки становления political культуры>>} \\
        \vspace{0.1cm}
        {\small по дисциплине <<Политология>>}
    \end{spacing}
\end{center}

\vspace{0.5cm}

Данный документ представляет собой статейный вариант реферата, отражающий ключевые тезисы, теоретико-методологические основания и выводы проведенного исследования, посвященного генезису и эволюции политической культуры.

\section{Понятие, структура и функции политической культуры}
В качестве концептуального базиса исследования принято системное определение А.И. Соловьева, согласно которому политическая культура трактуется как относительно устойчивая, исторически сложившаяся система ценностей, убеждений, традиций и паттернов поведения. Выделена трехкомпонентная внутренняя структура феномена:
\begin{enumerate}
    \item \textbf{Познавательно-нормативный компонент} (политические знания, правовые нормы, ценностные идеалы);
    \item \textbf{Оценочно-эмоциональный компонент} (субъективное отношение граждан к институтам власти, лидерам и государственной символике);
    \item \textbf{Практико-поведенческий компонент} (устойчивые стереотипы участия, электоральная активность, формы взаимодействия с государством).
\end{enumerate}
Определена функциональная матрица феномена, включающая функции гражданской идентификации, преемственной социализации, нормативной регуляции и системной интеграции общества.

\section{Классические теоретические подходы (Алмонд, Верба, Хантингтон)}
На основе компаративного анализа классических политологических концепций сформулированы два фундаментальных тезиса:
\begin{enumerate}
    \item \textbf{Концепция <<гражданской культуры>> (Г. Алмонд, С. Верба):} доказано, что стабильность демократии не обеспечивается <<чистыми>> типами культур (патриархальной, подданнической или активистской). Оптимальным признан смешанный тип --- \textit{гражданская культура} (\textit{Civic Culture}), гармонично балансирующая активное субъектное участие масс с подданнической лояльностью государственным институтам.
    \item \textbf{Концепция политического порядка (С. Хантингтон):} в условиях глобальной модернизации стабильность макросистемы детерминирована соотношением между мобибилизацией населения и уровнем институционализации. Разрыв между ростом активности масс и слабостью легитимных каналов ее интеграции ведет к деградации порядка и переходу от \textit{гражданских систем} к нестабильным \textit{преторианским системам}, основанным на прямом неправовом действии.
\end{enumerate}

\section{Историко-философская эволюция политических ценностей}
Генезис аксиологической матрицы гражданственности рассмотрен в диахроническом срезе:
\begin{enumerate}
    \item \textbf{Античность (Платон, Аристотель):} заложены истоки нормативно-активистской модели, провозгласившей человека <<политическим животным>>, а участие в делах полиса --- высшей этической обязанностью.
    \item \textbf{Средневековье (А. Августин, Ф. Аквинский):} доминирование теоцентрического подданничества, легитимация конформизма и покорности власти как божественному установлению.
    \item \textbf{Возрождение (Н. Макиавелли):} секуляризация политической мысли, переход к прагматической, рационально-технологической модели элитарного поведения.
    \item \textbf{Новое время (Т. Гоббс, Дж. Локк, Ж.-Ж. Руссо):} формирование концепции общественного договора, переход от подданничества к гражданственности, фиксация приоритета естественных прав личности над государством.
\end{enumerate}

\section{Факторы и институты политической социализации}
Воспроизводство и трансляция политико-культурных кодов осуществляются через систему агентов социализации. \textbf{Семья} выступает первичным латентным фактором, закладывающим устойчивые психоэмоциональные установки отношения к авторитету. \textbf{Система образования} целенаправленно формирует когнитивный (познавательный) пласт культуры через трансляцию исторической памяти. В современном обществе доминирующим агентом становятся \textbf{СМИ и цифровые коммуникации}, осуществляющие медийное конструирование виртуальной реальности и фрейминг политической повестки дня.

\section{Характерные черты и трансформация политической культуры России}
На основе историко-генетического анализа верифицированы устойчивые детерминанты ядра отечественной политической традиции:
\begin{enumerate}
    \item \textbf{Этатизм (государствоцентризм)} --- восприятие государства как сверхидеи, высшего арбитра и патерналистского гаранта порядка в ущерб индивидуальной автономии граждан;
    \item \textbf{Персонификация власти} --- привязка легитимности институтов к конкретной личности лидера, а не к формально-правовым процедурам;
    \item \textbf{Фрагментированность (расколотость)} --- ценностный дуализм между модернизационными ориентациями элит и глубоким традиционализмом масс;
    \item \textbf{Инверсионный характер участия} --- сочетание длительных периодов пассивной лояльности со вспышками радикального неинституционализированного протеста.
\end{enumerate}
\textbf{Вывод.} Современный этап эволюции политической культуры России характеризуется сложным транзитным синтезом: традиционные патерналистские архетипы постепенно дополняются рационально-субъектными, активистскими ориентациями (рост локальных гражданских инициатив, цифровой активизм), что предопределяет вектор долгосрочной модернизации всей политической системы страны.

\end{document}